% ---------------------------------------------------------------------------
\section{Comparison with people who actually wrote by hand}
\label{sec:scribes}

Many of the properties above could, in principle, be normal habits of a scribe: people writing by hand abbreviate more
at the end of a line, alternate letter forms, change habits from one day to the next. A systematic comparison with real
scribes was missing, because it requires transcriptions that preserve letter forms and the original lines and pages. We
found them in two public archives: Menota (Nordic manuscripts) and the ReF (German manuscripts). To these we added the
copyist of the Copiale cipher, who chooses between equivalent symbols as someone writing in a cipher might.

For each property we used \textbf{the same measure} as for the Voynich, recomputing the Voynich in the same run.

\begin{table}[ht]
\centering\footnotesize
\caption{The ``scribe battery'': properties of the Voynich measured identically on real scribes and on the Copiale
copyist.}
\label{tab:battery}
\begin{tabularx}{\linewidth}{>{\raggedright\arraybackslash}p{2.9cm}YYYYY}
\toprule
property & Voynich & 6 Nordic scribes (1200--1550) & 73 German scribes (1350--1500) & Copiale copyist & verdict \\
\midrule
junction rule (effect / uncertainty of the choice) & 0.058 & 2 of 11 choices at 0.013--0.015, others zero & 2 of 13 at
0.001--0.005, others zero & --- & \textbf{proper to the Voynich} (4--12 times the strongest scribe) \\
left margin: lines avoiding the start of the line above & 5 initials avoided & 0 of 72 tests & no manuscript with
$\ge 3$ initials avoided; 3 of 72 with one & --- & \textbf{proper to the Voynich} \\
line-edge forms (rise of the most typical form) & end \eva{-m} $+0.151$; start \eva{s-} $+0.098$ & $+0.016$--$0.025$
(abbreviations at line end, \eva{v-} at start) & $+0.019$ (line-break hyphen); $+0.012$ & --- & same nature, \textbf{4--8
times stronger} in the Voynich \\
repetition of the adjacent word (identical / near) & 1.00\% / 3.87\% & 0.01--0.05\% / 0.2--0.5\% & median 0.05\% /
0.46\% & 0.00\% / 0.04\% & \textbf{proper to the Voynich} \\
drift along the line (per 10 glyphs) & $-0.020$ to $-0.034$ & at most 0.0075; abbreviation $+0.005$ & at most 0.008,
in no common direction & --- & \textbf{proper to the Voynich} \\
short state of choices (adjacent; 2--3 words) & $+0.13$; $+0.09$, equal for different words & 1 of 18 choices with a
similar window, but coming from similar words & no window in 9 choices (3 with few data); \eva{í/i} agrees only between
adjacent words & \textbf{rotates} homophones: $-0.18$ & \textbf{not found} in
scribes in the Voynich's form \\
slow state between consecutive lines & $+0.050$ & 6 of 17 choices agree across lines, 5 as large as the Voynich & --- &
--- & \textbf{normal for handwriting} \\
\bottomrule
\end{tabularx}
\end{table}

\subsection{What it means}

\begin{itemize}
\item \textbf{Five properties out of seven have no counterpart} in 79 scribes of three centuries and four traditions
  (Icelandic, Norwegian, Swedish, German), nor in the copyist of a cipher: junction rules, left margin, repetitions and
  drift along the line; line-edge forms have the same nature (a
  real scribe abbreviates more at line end) but are four to eight times weaker.
\item \textbf{The Voynich's short state} (agreement that holds even between very different words) was found in none of
  the 79 scribes. In the German manuscripts, measured together, none of nine letter choices has it (for three rare
  choices the data are too few to tell); one choice (\eva{í/i}) shows agreement between adjacent words only, without the
  two-to-three-word persistence of the Voynich.
  One Norwegian scribe (AM~302~fol, around 1300) has, in the choice between long and round \emph{s}, a ``window'' of
  similar size and shape, which passes the same tests as the Voynich (it does not come from short stretches of the page
  nor from position); but his window comes mainly from \textbf{similar words} ($+0.29$ between similar words, $+0.065$
  between different words), that is, from repeating forms of the same word: a different mechanism.
\item \textbf{The Copiale copyist does the opposite of the Voynich:} if in the previous word he used one symbol for a
  letter, in the next word he tends to use another (agreement $-0.18$ between adjacent words, down to $-0.48$ for
  \emph{u}). It is the practice of someone who encrypts with homophones and rotates them so as not to leave repetitions.
  The Voynich choices do not behave like the homophones of this cipher.
\item \textbf{The slow state} that carries from one line to the next is instead \textbf{normal} for handwriting (habits
  that change slowly, or changes of hand not marked in the files). It does not distinguish the Voynich.
\item \textbf{Repetitions do not come from genre either:} fifteenth-century German herbals, recipe books and astrology
  texts written by hand repeat the adjacent word 7 to 30 times less than the Voynich.
\end{itemize}

\noindent In one sentence: \textbf{the ``way of writing'' of the Voynich is not that of a scribe copying a text in a
language with ordinary variant forms.} Its rules between neighbouring words and at line edges are much stronger. This
does not say what the text is; it says that whoever wrote it followed rules that ordinary scribes did not follow.

% ---------------------------------------------------------------------------
\section{What the Voynich is not}

Every ``no'' in this section holds because the method used says ``yes'' when the answer is there (positive control).
Where a control did not work, the test was not used.

\subsection{Not a language written in an ordinary way}

\begin{itemize}
\item The statistical fingerprint (Section~4.1) matches none of the roughly one hundred languages of the sample, nor
  medieval technical texts.
\item \textbf{The space is not lexical} (Section~4.2): in languages the space depends on the word, here on the
  neighbouring glyphs.
\item \textbf{Immediate repetition and its shape} (Section~6.1): languages avoid repeating at once, the Voynich repeats at
  once.
\item \textbf{No link beyond the adjacent word} (Section~4.3), whereas languages have links at a distance as well.
\item \textbf{An honest caveat:} with the corrected measure, the agreement between Voynich choices has the same size as
  grammatical agreement in languages (for example between the endings \emph{-o/-a} in Italian or \emph{-us/-a} in Latin),
  and its shape (how long it lasts at two-three words) lies at the top of the range of languages, not outside it: the
  technical Latin of Pliny comes close. The short state alone is not enough to distinguish the Voynich from a language
  with agreement; the other properties together are.
\end{itemize}

\subsection{Not a real text encoded in a simple way}

We encoded real texts (Latin, Italian, the Bible, technical texts) in many ways and measured whether they acquire the
properties of the Voynich:
\begin{itemize}
\item \textbf{word-for-word code} (each word replaced by a Voynich word): it takes the letters and vocabulary of the
  Voynich, but not the repetitions, page homogeneity or junction rules;
\item \textbf{verbose ciphers} (each letter becomes a group of 2--3 glyphs): of 450 variants none reaches the
  predictability of the Voynich with words of its length;
\item \textbf{the Naibbe cipher}: similar in predictability and word length, not in repetitions, unique words, page
  homogeneity, junction (almost zero) and line properties;
\item \textbf{abbreviations, null letters, transpositions}: they move away or do not suffice;
\item \textbf{lists, prayers, litanies} encoded: they take on the ``grain'' of the Voynich but not its anomalies.
\end{itemize}

\subsection{Not a substitution readable in a known language}

\begin{itemize}
\item \textbf{Homophonic substitution with spaces in 14 languages:} in the controls (the same kind of cipher on another
  biblical text) 69--98\% of real words are recovered; on the Voynich at most 36\%, with 7--93 distinct words, mostly
  always the same three, as for a text in another language.
\item \textbf{A solver without spaces} (simulated annealing with a 5-gram model of Latin trained on 9 million letters),
  calibrated on 56 + 71 positive controls (it recovers the key at 98--100\%): on the Voynich, counting glyphs in four
  different ways and at nine levels of grouping, in 71 languages, even reading backwards, the result never reliably
  exceeds that of the negative controls. The ``deciphered'' text is a mush of syllables.
\item \textbf{Other tests:} 85 languages with spaces; 16 vernaculars, historical languages and abbreviated Latin; a mixed
  nomenclator; searches guided by the expected content (plant names, months); other reading orders; pieces of words as
  symbols: no reading. Two ``formally positive'' outcomes (a German candidate and a syllabic one) proved to be artefacts:
  the second came out identical on the shuffled Voynich.
\item \textbf{Ordered anagrams:} the order of glyphs inside words is language-like, not anagram-like.
\item \textbf{Plant names as cribs:} the name of the plant, enciphered on purpose in a control, is recovered in the first
  word of the page; in the Voynich, with the available identifications, there is no trace.
\end{itemize}

\subsection{The published readings we could test}

\begin{itemize}
\item \textbf{Bax (2014):} the recurring words do not behave as names (\eva{shor}, read as ``black'', occurs 96 times on
  67 pages of all sections); outside the pages it was derived from, the key does not beat chance ($p = 0.14$) and
  produces ``names'' everywhere.
\item \textbf{Vatne (2021):} it relies on its own transcription; only a third of its entries are found in the reference
  transliteration.
\item \textbf{Cheshire (2019, ``proto-Romance'' language):} his key reads more Romance words than chance, but \emph{any}
  key of the same form adapted to the manuscript in a few minutes does better, even on pages not used to adapt it (20 keys
  out of 20).
\item \textbf{Schechter (2026, glossary from EVA to Latin, Occitan and Hebrew;} according to the author a bilingual
  Latin--Occitan pharmaceutical manual, ``87.8\% decoded''): his program, re-implemented, gives the same figures; but a
  ``glossary'' made of the most frequent words without any meaning covers as much, his glossary ``reads'' 67--70\% of a
  text without a message as well, and the order of the decoded words is not Latin (attested pairs as in shuffled text).
\item \textbf{Gatta (2026, a mapping of 19 EVA glyphs to Hebrew consonants, read right to left):} the signal present on
  the Voynich is present on the text without a message too, and a mapping searched for on purpose does better than his.
  The author himself, in the current version of his work, concludes that the text does not behave like Hebrew and that no
  page is readable.
\end{itemize}
The \textbf{message-free text control} (applying the same reading to a text generated without content) is simple, and we
propose it as a minimum test for anyone announcing a decipherment.

\subsection{None of the simulated historical cipher mechanisms}

We looked in the Voynich for the signature that some fifteenth- and sixteenth-century encryption systems would leave,
after checking on texts enciphered on purpose that the test can see it:
\begin{itemize}
\item \textbf{homophones chosen by the copyist} (as in the Copiale): they would leave a \emph{negative} agreement between
  nearby words (the copyist alternates the symbols); the Voynich has a positive agreement;
\item \textbf{a sign that changes the alphabet} (the index letters of Leon Battista Alberti's cipher disk, \emph{De
  componendis cifris}, c.~1466): no frequent word or sign after which the state of the choices restarts (a rare index
  cannot be excluded);
\item \textbf{a bit message in the binary choices} (like Francis Bacon's biliteral cipher, devised c.~1576--79 and
  described in 1623, where the form of the letter carries a bit): a message restarting at each page or line would be
  seen with overwhelming evidence in the controls ($z$ up to 70); in the Voynich there is no signal. A continuous message
  crossing pages remains outside the reach of the test;
\item \textbf{word tables in turn} (the ``Ave Maria'' of Johannes Trithemius' \emph{Polygraphia}, printed in 1518, where
  each letter becomes a Latin word taken from the next table and the ciphertext reads like a prayer): no rhythm from two
  to six words along the line;
\item \textbf{an autokey restarting at each line} (each encrypted letter depends on the previous one; each line has a new
  starting key): it is the only mechanism tested that \textbf{closes the line} like the Voynich, but it destroys
  everything else: no repetition, an inflated vocabulary (4{,}500 to 6{,}400 distinct words per 10{,}000, against 2{,}900),
  weak line starts, no avoidance at the margin.
\end{itemize}

\subsection{Not a constructed language, not gibberish, not the published generators}

\begin{itemize}
\item \textbf{Constructed languages} (Esperanto, Volapük, Neo-Quenya, Interlingua, Lojban, Klingon, LOLCat): none has the
  state of the choices; none has even an agreement between adjacent words.
\item \textbf{Meaningless text handwritten by volunteers:} it has no junction (0.014), no line-edge forms, no copying from
  the line above, no avoidance of repeated starts (it repeats them). Someone inventing a text by hand does not
  spontaneously produce the rules of the Voynich.
\item \textbf{Published generators:} Naibbe has neither junction nor line properties; the Timm and Schinner generator has
  copying from the lines above and a non-lexical space, but it draws from too far, has no junction closed within the line,
  no state of the choices (its agreement comes from copying similar words), no margin avoidance and no drift; generators
  U2 and U3 lack the combination of chain and closed line. \textbf{None reproduces together} a junction closed within the
  line, short memory of repetitions, the state of the choices and margin avoidance.
\end{itemize}

% ---------------------------------------------------------------------------
\section{What we confirm and what we refute in previous work}

\begin{center}\footnotesize
\begin{tabularx}{\linewidth}{>{\raggedright\arraybackslash}p{5.2cm}Y}
\toprule
previous idea or result & our outcome \\
\midrule
\textbf{Currier's languages A and B} & \textbf{Confirmed} (98.5\% of pages; confirmed by the vocabulary in 107 of 107 herbal
pages). They differ in vocabulary, not in glyphs; each bifolium is entirely A or entirely B. \\
\textbf{Several hands} (Currier; Davis 2020) & \textbf{Not questioned}, but the fine habits distinguishing hands are partly
confounded with sections (hands 2 and 3 cannot be distinguished within the herbal). Above all, \textbf{all follow the same
rules} (junction, state of the choices, margin). \\
\textbf{Self-citation} (Timm and Schinner 2020): the scribe copies and modifies words already written & \textbf{Partly
confirmed.} The scribe does copy from the line above (more than their generator) and repeats at once. But Voynich reuse is
\textbf{short-range} (the line immediately above; in-line memory of 6--7 words), \textbf{adapted to the junction rules},
with a state of the choices that holds between different words: things their generator does not do. \\
\textbf{The Naibbe cipher} (Greshko 2025) as an explanation & \textbf{Not sufficient as a complete explanation}: no
junction, no line properties, repetitions and unique words far off. \\
\textbf{Topic ``keywords''} (Montemurro and Zanette 2013) & \textbf{Not distinctive}: texts generated without a message have
them too. \\
\textbf{Gibberish resembles the Voynich} (Gaskell and Bowern 2022, on low-level statistics) & \textbf{Not in the rules of
writing}: handwritten gibberish has no junction, no line-edge forms, no copying from the line above, no margin
avoidance, nor the regularities of the vocabulary (frequent words are not the most probable forms). \\
\textbf{``A glyph is not a letter, a token is not a word, a space is not a space''} (Rozanova and Temerev 2026) &
\textbf{Convergent}: we too find a chain of glyphs with rules at word edges, glyphs too regular for a one-to-one
substitution, and uncertain spaces that behave like word-internal junctures. Like them, we find that the ``words'' form a
plausible vocabulary and that the anomaly lies in how they follow one another. \\
\textbf{Readings by Bax, Vatne, Cheshire, Schechter, Gatta} & \textbf{Do not survive} the controls (Section~8.4). \\
\textbf{Zodiac labels as numbers or days} & \textbf{No}: almost all different, no positional correspondence between
months. \\
\textbf{``The Voynich has little syntax''} & \textbf{Not an anomaly}: Latin technical texts have as little. \\
\bottomrule
\end{tabularx}
\end{center}
